\chapter{Boundary-Value Problems for ODEs}

\section{Two-Point Boundary-Value Problems}

In a second-order initial-value problem, both $y$ and $y'$ are specified at the same point.  In a two-point boundary-value problem (BVP), values of $y$ are specified at different endpoints.

\begin{definition}[Two-point BVP]
A two-point boundary-value problem has the form
\[
 y''=f(x,y,y'),
 \qquad a\le x\le b,
\]
subject to
\[
 y(a)=\alpha,
 \qquad
 y(b)=\beta.
\]
\end{definition}

\begin{definition}[Linear second-order equation]
The equation is linear if it can be written as
\[
 y''=p(x)y'+q(x)y+r(x),
\]
or equivalently
\[
 y''-p(x)y'-q(x)y=r(x).
\]
\end{definition}

The two principal numerical approaches in this topic are the shooting method and the finite-difference method.

\section{Linear Shooting Method}

Consider
\[
 y''=p(x)y'+q(x)y+r(x),
 \qquad
 y(a)=\alpha,
 \qquad
 y(b)=\beta.
\]
The linear shooting method replaces the BVP by two IVPs.

First solve
\[
 y_1''=p(x)y_1'+q(x)y_1+r(x),
 \qquad
 y_1(a)=\alpha,
 \qquad
 y_1'(a)=0.
\]
Next solve the homogeneous IVP
\[
 y_2''=p(x)y_2'+q(x)y_2,
 \qquad
 y_2(a)=0,
 \qquad
 y_2'(a)=1.
\]
Because of linearity, every solution satisfying the left boundary condition has the form
\[
 y(x)=y_1(x)+cy_2(x).
\]
The right boundary condition gives
\[
 c=\frac{\beta-y_1(b)}{y_2(b)},
\]
provided $y_2(b)\ne0$.  Therefore
\[
 \boxed{
 y(x)=y_1(x)+
 \frac{\beta-y_1(b)}{y_2(b)}y_2(x)
 }.
\]
The two IVPs may be integrated by RK4 or another IVP solver.

\begin{example}
Apply linear shooting with RK4 and $h=1/3$ to
\[
 y''=4(y-x),
 \qquad0\le x\le1,
 \qquad y(0)=0,
 \qquad y(1)=2.
\]
Estimate $y(1/3)$ and $y(2/3)$.
\end{example}

\begin{solution}[9.0cm]
The first IVP is
\[
 y_1''=4(y_1-x),
 \qquad y_1(0)=0,
 \qquad y_1'(0)=0.
\]
RK4 gives
\[
\begin{array}{c|rr}
\toprule
x&y_1&y_1'\\
\midrule
0&0&0\\
1/3&-0.024691&-0.230453\\
2/3&-0.214398&-1.026741\\
1&-0.809732&-2.755580\\
\bottomrule
\end{array}
\]
The homogeneous IVP is
\[
 y_2''=4y_2,
 \qquad y_2(0)=0,
 \qquad y_2'(0)=1.
\]
RK4 gives
\[
\begin{array}{c|rr}
\toprule
x&y_2&y_2'\\
\midrule
0&0&1\\
1/3&0.358025&1.230453\\
2/3&0.881065&2.026741\\
1&1.809732&3.755580\\
\bottomrule
\end{array}
\]
The shooting coefficient is
\[
 c=\frac{2-y_1(1)}{y_2(1)}
 =\frac{2.809732}{1.809732}
 \approx1.552568.
\]
Therefore
\[
 y\left(\frac13\right)
 \approx-0.024691+1.552568(0.358025)
 =0.5312,
\]
\[
 y\left(\frac23\right)
 \approx-0.214398+1.552568(0.881065)
 =1.1535.
\]
\end{solution}

\section{Shooting for Nonlinear Problems}

For
\[
 y''=f(x,y,y'),
 \qquad y(a)=\alpha,
 \qquad y(b)=\beta,
\]
treat the unknown initial slope as a parameter:
\[
 y'(a)=z.
\]
For a chosen $z$, solve the IVP and denote the computed endpoint value by
\[
 y(b)=\phi(z).
\]
The correct slope solves the scalar nonlinear equation
\[
 F(z)=\phi(z)-\beta=0.
\]
Thus a nonlinear shooting method combines an IVP solver with a root-finding method such as bisection, secant or Newton iteration.  Each root iteration requires one or more integrations of the IVP.

\begin{algorithm}[Nonlinear shooting]
\begin{enumerate}
  \item Choose a trial slope $z$.
  \item Solve
  \[
    y''=f(x,y,y'),
    \qquad y(a)=\alpha,
    \qquad y'(a)=z.
  \]
  \item Compute the boundary residual $F(z)=y(b;z)-\beta$.
  \item Adjust $z$ with a root-finding method and repeat until $|F(z)|$ is sufficiently small.
\end{enumerate}
\end{algorithm}

\section{Finite-Difference Method}

Consider the linear BVP
\[
 y''+P(x)y'+Q(x)y=f(x),
 \qquad y(a)=\alpha,
 \qquad y(b)=\beta.
\]
Let
\[
 x_i=a+ih,
 \qquad h=\frac{b-a}{n},
 \qquad y_i\approx y(x_i).
\]
At each interior node $i=1,\ldots,n-1$, use central differences:
\[
 y''(x_i)\approx
 \frac{y_{i+1}-2y_i+y_{i-1}}{h^2},
\]
\[
 y'(x_i)\approx
 \frac{y_{i+1}-y_{i-1}}{2h}.
\]
Substitution gives
\[
 \frac{y_{i+1}-2y_i+y_{i-1}}{h^2}
 +P_i\frac{y_{i+1}-y_{i-1}}{2h}
 +Q_i y_i=f_i.
\]
After multiplying by $h^2$,
\[
 \left(1-\frac h2P_i\right)y_{i-1}
 +(h^2Q_i-2)y_i
 +\left(1+\frac h2P_i\right)y_{i+1}
 =h^2f_i.
\]
The boundary values $y_0=\alpha$ and $y_n=\beta$ are moved to the right-hand side.  The unknown interior values satisfy a tridiagonal linear system.

\begin{example}
Use the finite-difference method with $h=1$ to approximate the solution of
\[
 y''-\left(1-\frac x5\right)y=x,
 \qquad y(1)=2,
 \qquad y(3)=-1.
\]
Find $y(2)$.
\end{example}

\begin{solution}[4.2cm]
There is one interior point, $x_1=2$, with unknown $y_1\approx y(2)$.  Here
\[
 P(x)=0,
 \qquad
 Q(x)=-\left(1-\frac x5\right),
 \qquad
 f(x)=x.
\]
At $x=2$,
\[
 Q_1=-\frac35,
 \qquad f_1=2.
\]
The finite-difference equation is
\[
 y_0+(Q_1-2)y_1+y_2=f_1,
\]
because $h=1$.  Substituting $y_0=2$ and $y_2=-1$ gives
\[
 2-\frac{13}{5}y_1-1=2.
\]
Thus
\[
 -\frac{13}{5}y_1=1,
 \qquad
 y_1=-\frac5{13}\approx-0.3846.
\]
Therefore
\[
 y(2)\approx-0.3846.
\]
\end{solution}

\begin{remark}
Shooting preserves the differential-equation viewpoint but can be sensitive when the endpoint value changes greatly with the initial slope.  Finite differences convert the entire BVP directly into a sparse linear or nonlinear algebraic system and are often more robust for difficult boundary-value problems.
\end{remark}
